\documentclass[12pt]{article}

\usepackage[utf8]{inputenc}
\usepackage[french]{babel}
\usepackage[paperwidth=20cm,paperheight=10cm,margin=5mm]{geometry}

\pagestyle{empty}
\setlength{\parindent}{0mm}
\setlength{\parskip}{1em}

\begin{document}

\section*{Séquence 1 (2'15'')}

Bonjour à tous, et bienvenu dans le module sur l'automatisation des
calculs avec les workflows.

\newpage

L'automatisation est la clé de la reproductibilité.

Un travail interactif n'est jamais reproductible.

Vous pouvez bien sûr tenter de prendre des notes soigneusement, mais
vous allez faire des erreurs.

Et même si vous ne faites pas d'erreur, vous ne pouvez jamais être sûr
de ne pas avoir fait d'erreur.

Vous savez exactement ce que vous avez fait seulement si vous pouvez
refaire complètement votre calcul, d'un bout à l'autre.

\newpage

L'automatisation des calculs est un des rôles principaux des ordinateurs.

On peut l'appliquer à plusieurs niveaux, et traditionnellement on
utilise des noms différents pour les logiciels qui s'appliquent à
chaque niveau.

Un programme est un logiciel qui définit des structures de données et
applique des algorithmes complexes de façon réutilisable.

Un script implémente des algorithmes simples pour résoudre un problème
spécifique, et n'introduit pas d'abstraction de données.

Le notebook, que nous avons présenté dans notre MOOC d'introduction,
est une variante du script qui rajoute des commentaires et des
résultats.

Enfin, le workflow, qui est le sujet de ce module, sert à coordonner
des programmes ou scripts multiples, exécutés sur des ordinateurs
multiples, pour traiter des jeux de données multiples.

\newpage

Dans ce module, je présente d'abord les situations dans lesquelles
les workflows se montrent utiles.

Par la suite, j'évoque les critères de choix d'un gestionnaire de
workflow, parce qu'il y en a un très grand nombre.

Je décrirai après la plus grande famille de gestionnaires de workflow,
dont l'ancêtre commun est un logiciel des années 1970s qui s'appelle
``make''.

Dans cette famille, il y a notamment ``snakemake'', qui est
particulièrement adapté au calcul scientifique.

J'appliquerai snakemake à l'analyse de l'incidence du syndrome
grippal, que nous avons déjà utilisé comme exemple dans notre premier
MOOC.

Enfin, je vais parler de l'intégration continue, qui est une forme
spécial de wirkflow pour tester un logiciel tout au long de son
évolution.

\newpage

\section*{Séquence 2 (2'35'')}

Dans cette séquence, je vous explique pourquoi les workflows ont été
inventés, et en quoi ils se distinguent des scripts et notebooks que
vous connaissez déjà.

\newpage

La difficulté qui se pose avec les scripts et notebooks et le passage
à l'échelle.

Quand un script devient trop long, il devient difficile à lire ou à
modifier, à cause de sa structure linéaire. Il faut sans cesse avancer
et reculer dans le code en cours de lecture.

Un workflow permet de découper le calcul en tâches plus petites, dont
chacune est implémentée dans un script de taille gérable.

\newpage

Une autre limite des scripts et notebooks est la durée d'exécution,
car il faut toujours les exécuter en entier. Prenons l'exemple d'un
calcul paramétrique qu'on applique à 100 valeurs pour le paramètre.
Puis on veut rajouter une 101ème valeur. Avec une boucle dans un
script, il faut refaire les 100 calculs déjà faits.

La même boucle dans un workflow est executée de façon plus intelligente :
quand le résultat pour une valeur donée du paramètre est déjà disponible,
le calcul n'est pas fait une deuxième fois.

\newpage

Les briques de calcul dans un workflow sont appelés des
``tâches''. Chaque tâche est un mini-calcul: un bout de code est
appliqué à des données d'entrée et produit des données de sortie. Une
tâche est typiquement réalisée par un script, et les données d'entrée
et sortie sont le plus souvent stockées dans des fichiers.

\newpage

Le workflow est alors le plan d'exécution des tâches, qui tient compte
du fait que l'entrée d'une tâche est la sortie d'une autre.

Un workflow ne permet que des structures de contrôle simples, par exemple
des boucles.

Un workflow peut être défini par un langage de programmation dédié,
mais aussi construit de façon interactive comme un graphe reliant des tâches.

\newpage

Voici un exemple d'un graphe, qui correspond au workflow que je
présenterai en détail plus tard. Ce workflow est définit par un langage,
le graphe ne sert que pour la visualisation.

Chaque boîte est une tâche, chaque ligne représente une dépendence
entre deux tâches. Le graphe montre bien la présence d'une boucle,
dont chaque itération est indépendente des autres, ce qui permet une
exécution en parallèle.

\newpage

Le gestionnaire de workflow qui est chargé de l'exécution du plan
commence par une analyse des dépendences, pour assurer qu'aucune tâche
n'est exécutée avant que ses données d'entrée ne soient disponibles.

Il décide quelles tâches doivent être exécutées, et dans quel ordre.

Il évite aussi d'exécuter des tâches dont le résultat est déjà disponible.

Il exécute les tâches en parallèle autant que possible, en fonction
de la disponibilité de moyens de calcul.

Enfin, il gère les erreurs de toute sorte qui peuvent se produire: une
panne de matériel, une coupure d'électricité, etc.

\newpage
\section*{Séquence 3 (5'00'')}

Dans cette séquence, je fais le tour des caractéristiques qui font la
différence entre les nombreux gestionnaires de workflow, afin de vous
aider à faire un bon choix.

\newpage

L'univers des workflows est en effet un peu bordélique. Il y a des
dizaines de gestionnaires populaires, et des centaines qui le sont
moins.

L'idée a été découverte de nombreuses fois, et beaucoup de gens ont
écrit leur gestionnaire personnel, spécifiquement adapté à leurs
besoins qui sont très variables d'un usage à l'autre.

Comme j'ai déjà dit, il n'y a aucune compatibilité de worklows entre
deux gestionnaires différents. Pour cette raison, vous n'avez
peut-être aucun choix vous-mêmes : si vous travaillez en équipe, vous
héritez souvent du choix de quelqu'un d'autre, parce que vous devez
travailler avec ses workflows.

Il y a une norme en cours de développement, le Common Workflow
Language, mais pour l'instant elle n'est pas très répandue.

\newpage

Un premier critère important dans le choix d'un gestionnaire
de workflows est la nature des ressource de calcul que vous
voulez mettre en oeuvre.

Si vos calculs peuvent tourner sur votre ordinateur personnel,
vous avez peu de contraintes car tous les gestionnaires proposent
ce mode de fonctionnement.

Mais si vous utilisez des ordinateurs à distance, par exemple dans
votre laboratoire, ou des ordinateurs dans une grille
ou dans le cloud, faites bien attention à ce que votre gestionnaire
sait parler à ces machines.

\newpage

Une question liée concerne l'ordinateur sur lequel tourne le
gestionnaire avec son interface utilisateur. C'est souvent votre
ordinateur personnel, mais quelques gestionnaires présentent une
interface Web et sont pilotés à partir d'un simple navigateur.

Ce choix vous facilite la gestion des ressources, parce qu'elle
devient le travail de quelqu'un d'autre, mais il pose aussi des
contraintes dans le choix des logiciels auxquels vous avec accès dans
vos workflows.

\newpage

La troisième question liées aux ressources porte sur le stockage des
données.

Pour des données de faible volume, ce n'est pas critique car vous
pouvez les copier là où il faut.

Par contre, le traitement de grands volumes de données nécessite
l'organisation des calculs autour du stockage des données, car le
transfert de données peut devenir plus couteux que le calcul lui-même.

En ce cas, vous avez besoin d'un gestionnaire de workflows qui sait
gérer les ressources dont vous disposez: des fichiers sur votre
ordinateur personnel ou sur un ordinateur de laboratoire, un dépôt
dans le cloud, genre AWS, ou des bases de données publiques
accessibles par le Web.

\newpage

Passons à une autre catégorie de critères: la description des calculs.

Une tâche est un petit calcul, qui correspond à l'exécution d'un
programme ou d'un script.

Les gestionnaires de workflows n'imposent pas eux-mêmes des
contraintes quant aux logiciels utilisés par effectuer les tâches,
mais c'est à vous d'assurer la disponibilités de ces logiciels et de
leurs dépendances sur les ordinateurs qui exécutent les tâches.

Certains gestionnaires vous aident dans cette gestion, par exemple en
gerant des conteneurs pour les environnements de calcul.

Pour la reproductibilité, cette gestion de l'environnement logiciel
de vos workflows est très importante.

Il y a aussi des gestionnaires de workflows qui acceptent des services
Web dans la définition d'une tâche, ce qui peut être un critère de
choix aussi.

Mais faites attention à la reproductibilité. Un service Web peut
disparaître d'un jour à l'autre, ou donner des résultats différents
pour les mêmes données d'entrée. En fait, le serveur qui propose le
service devient un élément de votre environnement logiciel!

\newpage

Quant à la définition du workflow lui-même, il y a trois façon de s'y
prendre.

La première est un langage spécifiquement fait pour décrire les workflows.
C'est le cas de snakemake, que j'utiliserai plus tard pour mon exemple,
vous verrez donc en détail ce mode de fonctionnement.

Une autre approche est l'utilisation d'un langage de programmation
généraliste, comme par exemple Python. Le gestionnaire de workflows
est alors une bibliothèque qui gère essentiellement les ressources de
calcul. Un langage généraliste donne plus de flexibilité, mais les
workflows deviennent plus longs et moins lisibles, et on fait plus
facilement des erreurs.

Enfin, les gestionnaires graphiques vous laissent dessiner directement
le graphe des dépendances entre les tâches. Comme souvent, une telle
approche visuelle facilite le démarrage, mais se montre plus
contraignante à terme. C'est notamment le cas dans un contexte collaboratif
qui nécessite une gestion de versions collaborative. Pour le moment,
cette technologie est bien développée uniquement pour les données
textuelles.

\newpage

Enfin, un critère plus technique est la flexibilité qu'un gestionnaire
de workflows vous donne pour définir l'algorithme global qui coordonne
vos tâches.

La gestion automatique des dépendances entre les tâches est un
minumum, tous les gestionnaires l'assurent.

Les boucles sur les valeurs d'un paramètre sont aussi presque
universellement disponibles.

Une fonctionnalité plus rare est la modification dynamique du graphe
des tâches. Ceci veut dire que vous pouvez rajouter ou supprimer une
tâche en fonction du résultat d'une tâche précédente. Je vous
montrerai plus tard dans mon exemple à quoi ça peut servir.

\end{document}
